\chapter{Object Reconstruction}
\label{chap:ObjectReco}

The raw data output of the ATLAS experiment consists of electrical signals from $\approx 40$ million channels. 
Particle trajectory and energy cluster information are combined in a Particle Flow (PFlow) algorithm to construct objects that serve as input to jet clustering algorithms.
Objects, such as jets, leptons and \ETmiss are reconstructed as high-level objects for
physics analyses. The objects are subsequently calibrated to account for differences
in MC simulation and data and to ensure accurate and precise measurements of their
properties.

This chapter details the reconstruction of tracks, calorimeter cluster and physics objects within the ATLAS detector.
The subsequent calibration will be discussed for the example of $b$ tag calibration in Chapter ...
The calibration of other objects will not be discussed
in detail.


along with the reconstruction and calibration
procedures for physics objects used in analysis.
%From alice:
%The identification of particles was previously introduced in Section 3.2.2; this section
%will focus on the reconstruction and calibration of physics objects relevant to the analyses
%presented in Chapters 5-7.

\section{Reconstruction of Tracks, Clusters and Vertices}
This section provides a concise summary of low level objects that will serve as an input for the reconstruction of physics objects.

\vspace{.5cm}
\noindent
\textbf{Tracks}
\vspace{.5cm}

\noindent
As illustrated in Figure~\ref{fig:ATLAS:DiagramParticles}, charged particles traversing
the various layers of the ATLAS detector produce spatial hits and energy depositions.
The particles are forced on helical trajectories by the
strong magnetic field, enabling the determination of momentum and charge.
The trajectories are reconstructed from ID and MS signals to form tracks.

The first step in the reconstruction of charged particles is the creation of three-dimensional \textit{space points} from neighboring cells in Pixel and SCT that exceed a certain threshold~\cite{PERF-2015-08}.
%The pixel clusters in the barrel are about two pixels in the $r$-$\varphi$ plane and from one to three pixels in the $\eta$ plane.

Sets of three space points form \textit{track seeds} which allow a first crude momentum estimate while maximizing the number of possible combinations.
The candidates are extrapolated by assuming a perfect helical
trajectory in a uniform magnetic field to estimate the distance to beam line. These are also called \textit{transverse impact parameter} $d_0$ and \textit{longitudinal impact parameter} $z_0$.

After applying a first set of momentum and impact parameter requirements, 
a combinatorial Kalman filter~\cite{FRUHWIRTH1987444} is used to
build track candidates from the chosen seeds in an iterative fitting process.
It incorporates additional space-points from the remaining
layers of the Pixel and SCT detectors which are compatible with the preliminary trajectory.
Track candidates created from random collection of space points are rejected to manage the available CPU budget for event reconstruction.

A \textit{track score} is defined evaluating tracks according to their momentum, number of clusters and $\chi^2$ of the fit, where more energetic tracks with an increased number of clusters are ranked higher.
The tracks a processed by an ambiguity solver in descending order of the track score to resolve overlapping space points.

The ambiguity solver rejects all tracks that fail a number of criteria
including kinematic cuts ($\pt > 400 \MeV$, $|\eta|<2.5$) and on the total 
number of clusters and holes. A hole is defined as an intersection of the reconstructed track trajectory with a sensitive detector element that does not contain a matching cluster. The detailed list can be found in~\cite{PERF-2015-08}.

The tracks can be extended in the outer TRT in the next step.
%A two-step approach is used.
An extension algorithm searches for seeded candidates before a second module
processes and evaluates these extensions.
It should be noted, that this is a pure extension, not a combination.
The silicon-only tracks are not modified~\cite{Cornelissen_2008}.

The first step in the reconstruction of muons in the MS~\cite{PERF-2015-10} is to form segments from hit patterns in the muon chambers.
A Hough transform~\cite{ILLINGWORTH198887} is used to search for aligned hits in MDT and trigger chambers.
It speeds up the reconstruction and makes it more robust against misidentification.

A straight line is fitted to the hits in each layer is performed to reconstruct the MDT segments.
Muon track candidates are then built by combining hits from segments in different layers. 
The segment-seeded combinatorial search starts from the middle
layers of the detector where more trigger hits are available,
and is then extended to use the segments from the outer and inner layers.
Finally, a global $\chi^2$ fit is performed to
the hits associated with each track candidate.

The information provided by ID, MS and calorimeters is combined, defining four types of muons~\cite{PERF-2015-10}, ranked by reconstruction quality:
\begin{itemize}
    \item Combined (CB) muons are formed with a global refit that uses the independent
    reconstructions in ID and MS.
    \item Segment-tagged (ST) muons classify an ID track as a muon if it is associated with
    at least one local track segment in the MDT or CSC.
    \item Calorimeter-tagged (CT) muons match a track in th ID with an energy deposit in the
    calorimeter compatible with a minimum-ionizing particle.
    \item Extrapolated (ME) muons are reconstructed from MS hits only, adding a loose
    requirement that the origin of the track is compatible with the IP.
\end{itemize}
Overlaps between the categories are resolved before producing the collection of muons used in physics analyses.

\vspace{.5cm}
\noindent
\textbf{Calorimeter cluster}
\vspace{.5cm}

\noindent

The reconstruction of particles depositing energy in the calorimeter system is 
based on a three-dimensional topological clustering of individual calorimeter cell signals~\cite{PERF-2014-07}.
The resulting topological cell clusters (\textit{topo-clusters}) carry both, shape and location information.

The cell significance $\varsigma^\text{EM}_\text{cell}$ is defined to control the 
cluster formation
\begin{equation}
    \varsigma^\text{EM}_\text{cell} = \frac{E^\text{EM}_\text{cell}}{\sigma^\text{EM}_\text{noise,cell}}
    \label{eq:ObjReco:topo_threshold}
\end{equation}
where $E^\text{EM}_\text{cell}$ and $\sigma^\text{EM}_\text{noise,cell}$
denote the cell signal and average (expected) noise, respectively.
The algorithm to form topo-clusters starts from seeds with highly significant signal
$|\varsigma^\text{EM}_\text{cell}|>4$ (primary seed threshold). Each of the seed cells forms a \textit{proto cluster}.
All neighboring cells with $|\varsigma^\text{EM}_\text{cell}|>2$ (growth threshold) are collected
iteratively, as well as cells with $|\varsigma^\text{EM}_\text{cell}|>0$ (principal cell filter).
In this case, neighboring cells refers to cells adjacent in
a given sampling layer, as well as cells with partial overlap in adjacent layers.

The resulting proto-cluster is characterized by a core of
seed cells, surrounded by an envelope of cells with less significant signals.
The energy of the massless cluster is given by the energy sum of the associated cells taking geometrical signal weights into account,
while $\eta$ and $\varphi$ are calculated from the energy-weighted average of the cells.

\vspace{.5cm}
\noindent
\textbf{Vertices}
\vspace{.5cm}

\noindent
As mentioned in Section ... , there are multiple proton-proton collisions per bunch crossing.
In order to identify the one collision of interest, the intersection of multiple tracks, referred to as \textit{vertex} is reconstructed~\cite{ATL-PHYS-PUB-2015-026}.
The vertices at the beam line are denoted as \textit{primary vertices}, where the one with the largest sum of squared transverse momenta of associated tracks is labeled the \textit{hard-scatter vertex}.

\textit{Secondary vertices} are displaced with respect to the beam line.
They can be an indication for long-lived particles such as $b$ hadrons, that travel a short distance in the detector before decaying.


\section{Electron Reconstruction}

\section{Muon Reconstruction}

\section{Jet Reconstruction}

\section{MET Reconstruction}


\begin{figure}[ht!]
    \centering
    \includegraphics[width=0.85\linewidth]{figures/ATLAS/cern_accelerator.pdf}
    \caption{Overview of CERN accelerator complex showing pre-accelerators and the LCH (dark blue) with the four major experiments ATLAS, CMS, ALICE and LHCb~\cite{Mobs:2636343}.}
    \label{fig:ObjectReco:Placeholder}
\end{figure}




\section{Definition of physics objects on particle level}

The aim of this section is to provide definitions of physics objects, such as leptons, jets and missing momentum, that allow for a direct comparison of theoretical and experimental results.
The ATLAS collaboration provides recommendations~\cite{ATL-PHYS-PUB-2015-013} for particle level definitions with a minimal model dependence.
They should be accurate, unambiguous and allow the comparison of predictions with measurement results between experiments without extensive prior experimental knowledge.
The specified particle level objects are as close to experimentally reconstructed objects as possible, 
avoiding large extrapolations to regions outside of the detector acceptance.

\begin{itemize}
    \item All \textbf{stable particles}, defined by a lifetime $c\tau_0 > 10$ mm within the angular acceptance $|\eta| < 5$ are considered.
    \item \textbf{Prompt leptons} are generated final-state neutrinos, electrons and muons not originating from hadron decays.
    Charged leptons from subsequent $\tau$ decays $(W\rightarrow \tau\nu_\tau \rightarrow \ell\nu_\ell\nu_\tau)$ are therefore also categorized as prompt.
    To take QED FSR into account, the charged leptons are \textit{dressed} by clustering all photons within a cone around the \textit{bare} lepton (after simulating QED FSR) with the anti-$k_t$ algorithm. The clustering distance parameter is set to $R=0.1$.
    \item \textbf{Jets} are build in a subsequent step by collecting all stable particles apart from the identified prompt dressed leptons and prompt photons. The anti-$k_t$ algorithm with $R=0.4$ is employed for the jet building.
    The jet is labeled as a $b$ jet if a $b$ hadron with $p_T>5\GeV$ is found within $\Delta R < 0.3$ around the jet axis. Otherwise it is labeled as a \textit{light} jet. No jet-lepton overlap removal is applied.
    \item \textbf{Missing transverse momentum} $\ETmiss$
    corresponds to the transverse vector sum of all neutrinos not from hadron decays.
\end{itemize}
